\section{Related works}
\label{sec:related}

\subsection{EEG-based emotion recognition}
EEG emotion recognition pipelines typically consist of artefact removal, segmentation, feature extraction, feature selection and classification \cite{Alarcao2019}. Frequency-domain features remain dominant. Duan et al.\ \cite{Duan2013} introduced differential entropy (DE) and reported 84.22\% accuracy against 76.56\% for the energy spectrum, and Zheng and Lu \cite{ZhengLu2015} used deep belief networks on DE features of the SEED dataset to identify critical frequency bands and channels, obtaining 86.08\% average accuracy with comparable performance from reduced 4--12 channel montages. Time--frequency decomposition with the DWT is a common alternative to Fourier analysis because it provides band-limited sub-signals with good temporal localisation; Murugappan et al.\ \cite{Murugappan2010} used Daubechies-4 wavelets to extract alpha, beta and gamma components and reached 83.26\% with a k-nearest-neighbour classifier on five discrete emotions. Nonlinear complexity measures such as sample entropy \cite{Richman2000}, permutation entropy \cite{Bandt2002}, symbolic dynamics entropy \cite{Zebrowski2000} and phase entropy \cite{Rohila2019} have been adopted because EEG is non-stationary and nonlinear. Jenke et al.\ \cite{Jenke2014} compared features from 33 studies and found that such advanced features outperformed plain band power and that multivariate feature selection performed slightly better than univariate selection. The mRMR criterion \cite{Peng2005} is a widely used multivariate filter that selects features with high mutual information with the class and low mutual information with one another.

A recurring concern in this literature is evaluation validity. Li et al.\ \cite{Li2021} demonstrated that block-designed EEG experiments, in which training and test samples originate from the same continuous recording block, allow classifiers to exploit temporal correlations unrelated to the stimulus class. For emotion recognition, a review of 101 DEAP studies reported that segmentation leakage, feature selection on the full dataset and related errors were common and could inflate accuracy substantially \cite{Khan2025}, and a review of 216 papers reported inconsistent splitting and metrics across the field \cite{Kukhilava2025}. Both reviews are preprints and are cited here as indicative evidence. The cross-session drop reported for SEED-IV, from 85.11\% within sessions to 72.39\% across sessions \cite{Zheng2019}, further illustrates how strongly the evaluation protocol affects reported performance.

\subsection{Eye tracking and emotion}
The pupil dilates in response to emotionally arousing pictures and sounds, both pleasant and unpleasant, and this dilation co-varies with skin conductance, indicating sympathetic activation \cite{Bradley2008,Partala2003}. Pupil size therefore indexes arousal rather than valence. It is also governed primarily by the light and near responses, of which the psychosensory component is comparatively small \cite{Mathot2018}. For dynamic stimuli such as films, luminance changes must be modelled or removed before pupil size can be interpreted as arousal \cite{Raiturkar2016}, and Lu et al.\ explicitly removed the light reflex before using pupil features for emotion recognition \cite{Lu2015}. Pupil and gaze data further suffer from blinks, data loss and limited precision, which must be measured and reported \cite{Holmqvist2012}; standard preprocessing guidelines address blink removal, interpolation and smoothing \cite{Kret2019}. Lim et al.\ \cite{Lim2020} surveyed eye-tracking emotion recognition and identified pupil diameter, fixation, saccade and blink features as the most commonly used, while noting that eye tracking alone remains a relatively immature approach.

\subsection{Multimodal EEG and eye-tracking emotion recognition}
Zheng et al.\ \cite{ZhengDong2014} combined EEG with pupil-diameter features from five subjects watching film clips and obtained 71.77\% (EEG), 58.90\% (eye tracking) and 73.59\% (feature-level fusion) for three classes. Soleymani et al.\ \cite{Soleymani2012mm} fused EEG with pupil response and gaze distance at decision level and, under a participant-independent protocol, reached 68.5\% and 76.4\% for three-class valence and arousal, respectively. Lu et al.\ \cite{Lu2015} compared feature-level, decision-level and fuzzy-integral fusion of 16 eye-movement features with EEG; the fuzzy integral, which combines classifier outputs according to the importance of each source \cite{ChoKim1995}, produced the best accuracy of 87.59\%. EmotionMeter \cite{Zheng2019} introduced the four-class SEED-IV dataset and showed that EEG and eye movements are complementary: EEG recognised happiness better, whereas eye movements recognised fear better. Liu et al.\ \cite{Liu2022} compared deep canonical correlation analysis and bimodal autoencoders across five datasets, adding weighted-sum and attention-based fusion, and achieved 87.5\% on SEED-IV. Fu et al.\ \cite{Fu2023} aligned 4-s EEG and eye-movement windows in a two-branch network with multi-scale attention fusion and reported 87.32\% on SEED-IV.

In all of these studies the eye-tracking signal is a source of \emph{features} or \emph{decisions} to be fused with EEG. The gated approach of Farhadi Sedehi et al.\ \cite{FarhadiSedehi2024a,FarhadiSedehi2024b} is different in kind: eye tracking determines \emph{which EEG segments} are analysed, so it acts as a temporal selector rather than as a feature source. Table~\ref{tab:related} summarises the reviewed multimodal studies. Accuracies are copied from the original papers and are not directly comparable, because datasets, class definitions and validation protocols differ.

\begin{table}[t]
\centering\scriptsize
\caption{Summary of related EEG and eye-tracking emotion-recognition studies (results as reported by the original authors)}\label{tab:related}
\begin{tabularx}{\linewidth}{@{}p{0.9cm}p{1.5cm}p{3.0cm}p{2.0cm}p{1.6cm}X@{}}
\toprule
Ref. & Dataset & Method & Use of eye tracking & Result & Limitation\\
\midrule
\cite{ZhengDong2014} & Own (5 subj., 3 cl.) & PSD/DE, SVM, FLF/DLF & Pupil features fused & 73.59\% & Very small sample$^\dagger$\\
\cite{Soleymani2012mm} & MAHNOB proto. (24 subj.) & EEG + pupil/gaze, DLF & Features fused & 68.5\% V, 76.4\% A & Trial-level only; no temporal selection$^\ast$\\
\cite{Lu2015} & SEED-style (9 subj., 3 cl.) & DE + 16 eye features, fuzzy integral & Features + fuzzy DLF & 87.59\% & Within-session; fusion at output level$^\ast$\\
\cite{Zheng2019} & SEED-IV (4 cl.) & DE/PSD, multimodal DNN & Features fused & 85.11\% & Cross-session drop to 72.39\%\\
\cite{Liu2022} & SEED, SEED-IV, SEED-V, DEAP, DREAMER & DCCA/BDAE, attention fusion & Features fused & 87.5\% (SEED-IV) & Learned modality weights, not per window$^\ast$\\
\cite{Fu2023} & SEED-IV & Two-branch net, multi-scale attention & Features fused & 87.32\% & Black-box; weights not temporal$^\ast$\\
\cite{FarhadiSedehi2024a} & SEED-IV, MAHNOB-HCI (4 cl.) & DWT entropies, mRMR, S-FIS-GA & Hard pupil gate selects EEG segments & 94.42\%, 85.20\% & Segment-level 5-fold CV; crisp gate$^\ast$\\
\cite{FarhadiSedehi2024b} & MAHNOB-HCI (2 cl.) & Granger causality images, ResNet-18 & Hard pupil gate & 91.00\% & Crisp gate$^\ast$\\
\bottomrule
\multicolumn{6}{@{}l}{\scriptsize $^\dagger$ stated by the authors; $^\ast$ our analysis. FLF/DLF: feature-/decision-level fusion; V/A: valence/arousal.}
\end{tabularx}
\end{table}

\subsection{Fuzzy logic in emotion recognition and soft gating}
Fuzzy sets \cite{Zadeh1965} represent graded membership, and Takagi--Sugeno--Kang (TSK) models \cite{TakagiSugeno1985} combine fuzzy antecedents with functional or constant consequents. Rules can be generated from data with table look-up schemes such as that of Wang and Mendel \cite{WangMendel1992}, which is close to the look-up-table rule extraction used in the baseline, and membership parameters can be learned with neuro-fuzzy systems such as ANFIS \cite{Jang1993}. In emotion recognition, fuzzy rules have been used as interpretable EEG classifiers \cite{Matiko2014,FarhadiSedehi2024a}, fuzzy linguistic summaries have been applied to EEG time series \cite{Kaya2024}, and fuzzy integrals have fused EEG and eye-movement classifiers \cite{Lu2015,ChoKim1995}.

Gating, the use of one signal to modulate the contribution of another, has been studied separately. The mixture-of-experts model \cite{Jacobs1991} uses a soft gating network to weight expert outputs, and gated multimodal units learn multiplicative gates that decide how much each modality contributes \cite{Arevalo2020}. Classical classifier-combination theory analyses weighted sum and product rules \cite{Kittler1998}, and quality-aware fusion incorporates measured signal quality into score combination \cite{Nandakumar2008}. Fuzzy support vector machines assign each training sample a fuzzy membership that scales its influence on learning, reducing the effect of noisy samples \cite{LinWang2002}. These works establish soft weighting as an alternative to hard selection, but their weights are either learned and opaque, applied to modalities rather than time windows, or derived from the training sample itself rather than from a second modality.

% =====================================================================
